
\section{Your Father is Merciful}

Luke 6:36

“Folkebladet,” July 18, 1917

There is a glorious consolation in this word of Jesus for you who have learned to know yourself and have discovered that you are a sinner, a great sinner, who by disobedience and wickedness has made yourself deserving of God’s wrath and punishment.

Luther gives himself and all of us a dreadful verdict in the Second Article: “lost and condemned to eternal death.” But is this verdict true? Can things really be so bad with us? Does the verdict fit us, fit you? Yes, it does! It is true what you and I have sung so many times:

“My sin is great and grievous quite,
from it my sorrows spring!”

But if the matter stands thus, that my sin, your sin against the living God, is so great and exposes you and me to his wrath and punishment, then this is a good word from Jesus: “Your Father is merciful.”

In our explanation, the question of what God’s mercy is receives this answer: “He has pity upon our misery.”

It is a beautiful answer, deep and true as far as it reaches. But it scarcely reaches far enough—that is, if we try to explain God’s pity from ordinary human relations. For human beings too can have pity upon one another and upon one another’s need and misery. But often it goes no further than “having pity,” or feeling sympathy and fellow-feeling with those whose need and misery we know. And why? Ah, partly because we stand powerless and helpless when the question becomes one of helping in the need, and partly because we are so often unwilling to make the sacrifice required of us if the need and misery of our fellow human beings, our brothers and sisters, is to be relieved and remedied.

With the Father, however, it is otherwise. He does not merely “have pity upon our misery”; he does not sit up there in his heaven and blessedness feeling sorry for us who are in misery and need, without being able and without being willing to remedy our deep misery. No, he is not like that. He possesses the power and the will by which he is able to save in the midst of misery. And he uses this power and this will.

“Then God in eternity had pity
upon this great distress of mine;
he remembered his mercy
and willed to provide me help.
He turned a fatherly mind toward me,
and then, without deceit, gave forth
the dearest thing that he possessed.”

“Your Father is merciful.”

Ah, you who long to find firm and secure ground upon which to build in life and in death—here it is: God’s mercy. “The eternal and incomprehensible One has unveiled the features of his face to us. This face is called Christ. With fathomless compassion Christ’s eternity-deep eyes look upon us. We know his compassion: ‘Though he was rich, for our sake he became poor, so that through his poverty we might become rich.’”

“O abyss! which all our sins
has swallowed and put to death through Jesus’ death,
and sends us only to repentance,
that we may find grace sweet.
Here Jesus’ bloody sweat cries out:
Mercy! Mercy!”

Yes, we surely think this word of Jesus is good: “Your Father is merciful.” Our hope, our salvation, our life are contained in it. Thanks be to God!

But—Jesus says something more: “Be merciful!” A command, an order from Jesus, then. And to whom? To all, all who know that they need “the Father’s” mercy. And commands, orders, we are not always fond of. Not even Jesus’ commands. Submit yourself, and you will see that it is so. But Jesus has said it, and he expects us at least to try to obey him.

What does he mean when he says: “Be merciful”? He means that we are to have pity upon the misery in which our brothers and sisters in the world are placed. And there is still so much misery, temporal and spiritual. And Jesus wills that this misery should go to our hearts, so that we are drawn by it and driven to do what we can to relieve it.

For he says: “Be merciful as your Father.”

Brothers and sisters, is our mercy akin to the Father’s? Ah, I am so afraid that many of us must lower our eyes in shame and answer, humiliating though it is: No.

Of the Father’s mercy we sing:

“The thought soon came to him
that the world must perish;
it cut him to the heart.
In such heartfelt love
he came down to us upon the earth
to relieve our pain.”

Such is his mercy. But ours? Does the thought cut us to the heart that many of our own people in this land live “without God and without hope”? Not only those who with deliberate intent have torn themselves away from God and chosen to live in rebellion against him, but the many who are deceived by sin and by the beguiling influence of the world, and who grieve over it and long for peace. It is surely true that there are many who go about with “secret weeping of the heart.” And here Jesus expects us to be merciful, to have pity upon them—yes, with such compassion as the Father’s, whose sacrifice was his only Son.

Yes, how far we still are from the goal! But let our recognition of this not tempt us to despondency, but rather to “press on to take hold of it because we have been taken hold of by Christ.”

Where, then, or in what way, can more of this self-sacrificing mercy become ours? I think I will answer: by looking long, long into the face of the merciful Father, deep, deep into the dying eyes of the rejected Savior, and remembering that it is he who says to us: “Be merciful!” I know it will help. Your heart will be enlarged to contain more of that compassion which drives you to fervent intercession also for the salvation of our own people, and drives you to send your gift to our home mission.

Be merciful—as your Father!

E. B.
